
\documentclass{article}

%--- Document Formatting ---%
\usepackage{lmodern} % gives you the more modern Latin Modern fonts
\usepackage[T1]{fontenc} % fontenc provides better and more character encoding
\usepackage[margin=3cm]{geometry} % geometry allows changing page sizes, orientation and margin
\usepackage{microtype}

\usepackage{titlesec} % titlesec allows changing the style, spacing and font of chapter, sections and subsections
% {left-indent}{space before}{space after}
\titlespacing*{\section}{0pt}{6ex plus 1ex minus .2ex}{2.3ex plus .2ex}
\titlespacing*{\subsection}{0pt}{2.5em plus 1ex minus .2ex}{2ex plus .2ex}
\linespread{1.2}


%--- Language ---%
\usepackage[australian]{babel}
\usepackage[useregional]{datetime2}


%--- Symbols ---%
\usepackage{amssymb} % amssymb provides symbols, mainly number systems
\usepackage{cancel} % adds \cancel{x}, allowing striking out derivations
%\usepackage{siunitx} % used for physics based units


%--- Tools and General Formatting ---%
\usepackage{mathtools} % an extension for amsmath, loads amsmath itself
\usepackage{amsthm} % improves proof based math
\usepackage{enumitem} % enumitem gives more control over lists

%\usepackage{graphicx} % use to include images/figures
%\usepackage{xcolor} % for colour
%\usepackage[hidelinks]{hyperref} % adds hyperlinks, can sometimes cause issues. Keep last
%\usepackage{cleveref} % extension for hyperref, auto formatted cross-reference



\begin{document}

\title{Proofs}
\author{Montague Hamilton}
\date{\today}
\maketitle

\section{Proofs}
\subsection{Proof by contrapositive}
Prove: if $n^{2}$ is even, then n is even \\
Let $n \in \mathbb{Z}$. Direct proof isn't possible, as it requires taking the square root of both sides of $n^{2} = 2k$. Contrapositive proof means if $n$ is odd, then $n^{2}$ must also be odd.
n = 2k + 1 is odd, because 2k is even, as any number divisible by 2 is even $\frac{2k}{2} = k$, and any even number plus one is odd. To prove $n^{2}$ is odd:
\begin{align*}
  n    &= 2k + 1 \\
  n^{2}   &= (2k + 1)^{2} \\
       &= 4k^{2} + 4k + 1 \\
  n^{2} &= 2(2k^{2} + 2k) + 1
\end{align*}
Theorem: If $n^{2}$ is even, then n is even. Proof(contrapositive): If $n$ is odd, then $n^{2}$ is odd. Let $n = 2k + 1$ for some integer k. Then $n^{2} = 4k^{2} + 4k + 1 = 2(2k^{2} + 2k) + 1$, which is odd. Hence $n^{2}$ forces $n$ even. $\blacksquare$ \\
\newline
Theorem: If $n^{2}$ is even, then $n$ is even. ($n^{2}$ even $\implies n$ even, $n \in \mathbb{Z}$) \\
Proof (contrapositive): If $n$ is odd, $n^{2}$ is odd. Let $n \in \mathbb{Z}$, then from the definition of an odd number, for some $k \in \mathbb{Z}$,
\[
  n = 2k + 1 
\]
Therefore,
\begin{align*}
  n^{2} &= (2k + 1)^{2} \\
     &= 4k^{2} + 4k + 1 \\
     &= 2(2k^{2} + 2k) + 1
\end{align*}
Since $(2k^{2} + 2k) \in \mathbb{Z}$, $n^{2}$ is odd. Thus the contrapositive is true, if $n^{2}$ is even, then $n$ is even.

Theorem: $n^{2}$ even $\implies n$ even, $n \in \mathbb{Z}, \therefore \neg n$ even $\implies \neg n^{2}$ even,
where $\neg n$ even $\iff n$ odd. $\therefore n$ odd $\implies n^{2}$ odd. $\blacksquare$


\newpage
\section{Questions}
\begin{enumerate}
  \item Spot the flaw: Claim $2 = 1$. Let $a = b$. Then $a^{2} = ab$ so $a^{2} - b^{2} = ab - b^{2}$, so $(a + b)(a - b) = b(a - b)$.
    Dividing both sides by $a - b: a + b = b$. Since $a = b: 2b = b$ so $2 = 1$ \\
    The obvious flaw in this logic is the division by $a - b$, which while $a = b$ means that we are dividing by 0.

  \item Theorem: If $n^{3}$ is even, $n$ is even. \\ 
    Proof (contrapositive): If $n$ is odd, then $n^{3}$ is odd. 
    Using the definition of an odd number, let $n = 2k + 1$, where $n,k \in \mathbb{Z}$. 
    \begin{align*}
      n^{3} &= (2k + 1)^{3} \\
         &= (2k + 1)(2k + 1)(2k + 1) \\
         &= 8k^{3} + 12k^{2} + 6k  + 1 \\
         &= 2(4k^{3} + 6k^{2} + 3k) + 1 
    \end{align*}
    Plugging 1 into the equation,
    \[
      2(4(1)^{3} + 6(1)^{2} + 3(1)) + 1 = 2(13) + 1 = 27
    \]
    Since $(4k^{3} + 6k^{2} + 3) \in \mathbb{Z}$, $n^{3}$ is odd. Thus the contrapositive is true, if $n^{3}$ is even, $n$ is even. $\blacksquare$

  \item Theorem: The sum of two odd integers is even. \\
    Proof (direct): Let $a, b \in 2 \mathbb{Z} + 1$
    By definition of odd numbers,  $a = 2k + 1$ and $b = 2m + 1$ for some $k, m \in \mathbb{Z}$
    \begin{align*}
      a + b &= (2k + 1) + (2m + 1)\\
            &= 2k + 2m + 2 \\
            &= 2(k + m + 1)
    \end{align*}
    $k, m \in \mathbb{Z} \implies k + m + 1 \in \mathbb{Z}$, so $2(k + m + 1)$ must be even $\therefore a + b$ is even, 
    so the sum of two odd numbers is even. $\blacksquare$

  \item Theorem: There is no largest even integer. \\
    Proof (contradiction): Suppose we have the largest $n \in 2 \mathbb{Z}$. 
    Since $n$ is even, $n = 2k$ for some $k \in \mathbb{Z}$. Consider the integer $n + 2$:
    \[
      n + 2 = 2k + 2 = 2(k + 1)
    \]
    Since $k + 1 \in \mathbb{Z}$, $n + 2 \in 2 \mathbb{Z}$. Furthermore, because $2 > 0$, we have $n + 2 > n$.
    This means that $n + 2 \in \mathbb{Z}$ and $n + 2 > n$, which contradicts our assumption that $n$ is the largest $\in \mathbb{Z}$.
    $\therefore$ There is no largest even integer. $\blacksquare$ 

  \item State the converse of $n > 2 \implies n^{2} > 4$ and decide if it is true - proof or counterexample. \\
    The converse is $n^{2} > 4 \implies n > 2$.\\
    Theorem: The statement $n^{2} > 4 \implies n > 2 $ is false. 
    Proof (counterexample): Let $n \in \mathbb{R}$, suppose $n = -3$ then $(-3)^{2} = 9$ which > 4, so $n^{2} > 4$.
    However, $-3 \le 2$, so $n \le 2$, $\therefore n^{2} > 4 \not \implies n > 2$. $\blacksquare$


\end{enumerate}

\section{Format testing}
\begin{enumerate}
  % Q2
  \item \textbf{Theorem:} If $n^{3}$ is even, $n$ is even

    \textbf{Proof (contrapositive):} If $n$ is odd, $n^{3}$ is odd. Let $n \in \mathbb{Z}$ By definition of odd numbers,
    \[
      n = 2k + 1 \quad\text{for some } k \in \mathbb{Z}
    \]
    So for $n^{3}$,
    \begin{align*}
      n^{3} &= (2k + 1)^{3} \\
            &= (2k + 1)(2k + 1)(2k + 1) \\
            &= 8k^{3} + 12k^{2} + 6k + 1 \\
            &= 2(4k^{3} + 6k^{2} + 3k) + 1
    \end{align*}
    Since $k \in \mathbb{Z}$ and integers are closed under addition and multiplication, 
    \[
      4k^{3} + 6k^{2} + 3k \in \mathbb{Z}
    \]
    Let $m = 4k^{3} + 6k^{2} + 3k$. Then $m \in \mathbb{Z}$ and
    \[
      n^{3} = 2m + 1
    \]
    By definition of an odd integer $n$ is odd. $n$ being odd $\implies n^{3}$ is odd. 
    This proves the contrapositive, $n^{3}\implies n$
    \hfill$\blacksquare$
    

  % Q3
  \item \textbf{Theorem:} The sum of two odd integers is even.

    \textbf{Proof (direct):}
    Let $a,b \in \mathbb{Z}$ be odd. By definition of odd numbers, 
    \[
      a = 2k + 1
      \quad\text{and}\quad
      b = 2m + 1
    \]
    for some $k,m \in \mathbb{Z}$

    \begin{align*}
      a + b
      &= (2k + 1) + (2m + 1) \\
      &= 2k + 2m + 2 \\
      &= 2(k + m + 1)
    \end{align*}

    Since $k,m \in \mathbb{Z}$ and integers are closed under addition, 
    \[
      k + m + 1 \in \mathbb{Z} 
    \]
     
    Let $n = k + m + 1$. Then $n \in \mathbb{Z}$ and
    \[
      a + b = 2n
    \]
    By definition of an even integer, $a + b$ is even. $\therefore$ the sum of two odd integers is always even.
    \hfill$\blacksquare$
   
  \item \textbf{Theorem:} There is no largest even integer.

    \textbf{Proof (contradiction):} Suppose we have the largest even integer $n$. Since $n$ is even,
    \[
      n = 2k \text{ for some } k \in \mathbb{Z}
    \]
    Consider the even integer $n + 2$,
    \[
      n + 2 = 2k + 2 = 2(k + 1)
    \]
    $k + 1 \in \mathbb{Z} \implies n + 2 \in \mathbb{Z}$. Furthermore, because $2 > 0$ then $n + 2 > n$. 
    
    So $n + 2$ is even, an integer and larger than $n$ which contradicts $n$ being the largest even integer.
    
    $\therefore$ there is no largest even integer. 
    \hfill$\blacksquare$

    % Q5
    \item \textbf{Theorem:} The statement $n^{2} > 4 \implies n > 2$ is false.

      \textbf{Proof (counterexample):} Let $n \in \mathbb{R}$. Suppose,
      \begin{gather*}
        n = -3 \\
        \text{then} \\
        n^{2} = (-3)^{2} = 9
      \end{gather*}
      $9 > 4$ so $n^{2} > 4$ is true. However, $-3 \le 2$, so $n \le 2$.
      \[
        \therefore n^{2} > 4 \not \implies n > 2 \qquad\blacksquare
      \]
\end{enumerate}

%testing
\end{document}
